\section{Measured bias dimensionality is a property of the contrast set}
\label{sec:geometry}

Section~\ref{sec:rank} found that intervention rank trades debiasing against quality. A
natural next step is to estimate, for a given behaviour, how many dimensions its
representation occupies, and to set the rank from that estimate, as our own pipeline
originally did.

The usual estimator is the singular value spectrum of a stack of contrastive
activation differences: collect pairs of prompts differing only in the
attribute of interest, take the difference of their activations, and read off
the top explained-variance ratio ($\mathrm{EVR}_1$) or the rank needed to
reach some fraction of the variance ($k@80\%$). Low $\mathrm{EVR}_1$ is read
as a diffuse behaviour needing many directions; high $\mathrm{EVR}_1$ as a
concentrated one that a single vector can address.

We ran four controls on this estimator. All four use the AccessEval
contrastive pool of 2{,}083 pairs on Llama-3.1-8B-Instruct, and all four
indicate that $\mathrm{EVR}_1$ and $k@80\%$ report properties of the
contrastive set rather than of the bias.

\begin{table}[t]
\centering
\small
\begin{tabular}{lrr}
\toprule
Contrast set ($n{=}13$, L21) & $\mathrm{EVR}_1$ & $k@80\%$ \\
\midrule
Heterogeneous (any query)      & $22.1\%$ & $7.0$ \\
Single category                & $22.3\%$ & $7.0$ \\
Single domain                  & $27.1\%$ & $5.9$ \\
5 base queries only            & $44.2\%$ & $3.1$ \\
2 base queries only            & $\mathbf{76.1\%}$ & $\mathbf{1.7}$ \\
\midrule
\multicolumn{3}{l}{\textit{Paper's reported Target A (L4+L15, $n{=}13$)}} \\
\quad reported profile         & $67.9\%$ & $2$ \\
\bottomrule
\end{tabular}
\caption{Holding the number of contrastive pairs fixed at 13 and varying only how many
distinct base queries they are drawn from. Means over 40 draws. The $95\%$ intervals are
$[17.8, 28.4]$ for the heterogeneous row and $[62.0, 87.5]$ for two base queries. The
profile the paper reports for DiscrimEval Target A falls inside the latter interval and
matches its rank. The effect replicates at L14, where the same contrast runs from
$22.9\%$ $[18.1, 28.6]$ to $69.4\%$ $[54.3, 80.4]$.}
\label{tab:homogeneity}
\end{table}

\paragraph{Prompt diversity moves the estimate more than anything else.}
We fix the number of pairs at $n{=}13$ and vary only how many distinct base
queries those pairs are drawn from (Table~\ref{tab:homogeneity}). Drawn from
the whole pool, $\mathrm{EVR}_1$ is $22.1\%$ with $k@80\%{=}7.0$. Restricted
to two base queries, the same benchmark at the same layer gives
$\mathrm{EVR}_1{=}76.1\%$ $[62\%, 87\%]$ with $k@80\%{=}1.7$. The effect
replicates at L14 ($22.9\%$ to $69.4\%$ $[54\%, 80\%]$). A benchmark
described as forty-dimensional produces a near-rank-1 spectrum when its
prompts are made homogeneous.

\paragraph{Three further controls agree.} Sampling 13 pairs by drawing whole base
queries rather than uniformly moves $\mathrm{EVR}_1$ from $21.5\%$ to $72.2\%$ on the
same data at the same layer. A null built by permuting the pairing, which destroys the
contrast while preserving both marginals, is \emph{higher} than the real data at
$n{=}13$ ($26.2\%$ against $21.5\%$) and separates from it only above $n \approx 250$,
so a concentration estimated from a dozen pairs is consistent with no contrastive
structure at all. And resampling pairs directly gives $k@80\%{=}40$ $[40, 41]$ where
resampling the 146 distinct questions gives $30.1$ $[27, 33]$, intervals that do not
overlap, so treating nested pairs as independent is not merely over-precise.
Appendix~\ref{app:controls} gives all three in full. The held-out test supplied a fifth
instance unprompted: the rank reaching $80\%$ of the variance is 33 and 36 at the two
layers when fit on the 633 held-out pairs, and 42 and 49 on the full 1{,}044, for the
same bias at the same layers.

\paragraph{What this means for the diagnostic.} Taken together, these
controls remove the basis for using $\mathrm{EVR}_1$ to prescribe an
intervention rank. Two benchmarks can differ in reported dimensionality
because one varies a single template while the other spans nine categories
and six domains, with no difference in the underlying bias. We had read the
same signal as evidence that bias geometry varies by bias type, and it does
not follow.

One earlier observation reads differently in this light. The first singular direction is
close to orthogonal to a trained probe for the same attribute ($\cos{=}{-}0.001$ at L21),
so the top of the spectrum is not measuring the bias either.
